% from ega3-6-fr.tex
\oldpage[III]{175}
propriété particulière, en dehors de la restriction précédente sur les
degrés. On peut alors se demander si inversement, étant donné un complexe
$\mathscr L_\bullet$ formé de $\mathscr O_Y$-Modules associés à des
$A$-modules projectifs de type fini, limité inférieurement et dont les termes
de degré $>0$ sont nuls, il existe un $Y$-schéma $X$, projectif et plat sur
$Y$ et un $\mathscr O_X$-Module localement libre $\mathscr F$, tels qu'il y
ait un isomorphisme $\mathscr H^\bullet(f,\mathscr F\otimes_Y\mathscr Q_\bullet)\xrightarrow{\sim}\mathscr H_\bullet(\mathscr L_\bullet\otimes_Y\mathscr Q_\bullet)$
fonctoriel en $\mathscr Q_\bullet$. L'intérêt d'un tel résultat serait de
réduire complètement la théorie cohomologique des Modules cohérents et
$Y$-plats sur des $Y$-schémas propres, à la théorie « à homotopie près » des
complexes de $A$-modules projectifs de type fini sur un anneau noethérien
$A$.
\end{remark}

% (continued in ega3-7-fr.tex, before its first page marker)
\section{Étude du changement de base dans les foncteurs homologiques covariants de Modules}
\label{section:III.7.fr}

\subsection{Foncteurs de $A$-modules.}
\label{subsection:III.7.1.fr}

\begin{env}[7.1.1]
Étant donné un anneau $A$ (non nécessairement commutatif), nous noterons
$\boldsymbol{Ab}_A$ la catégorie des $A$-\emph{modules à gauche}, et nous noterons
simplement $\boldsymbol{Ab}$ la catégorie des $\mathbf Z$-modules, identiques aux
groupes commutatifs. Soit $T:\boldsymbol{Ab}_A\longrightarrow\boldsymbol{Ab}$
un \emph{foncteur covariant additif}, et soit $M$ un $(A,A)$-\emph{bimodule}; $T(M)$ est
alors muni de façon naturelle d'une structure de $A$-\emph{module à droite}. En
effet, pour tout $a\in A$, notons $h_{a,M}$ (ou simplement $h_a$)
l'endomorphisme $x\to xa$ du $A$-module \emph{à gauche} $M$. Par hypothèse,
$T(h_a)$ est un endomorphisme du $\mathbf Z$-\emph{module} $T(M)$; en outre, comme
$T$ est un foncteur covariant additif, on a, pour $a\in A$, $b\in A$,
\[
  T(h_{ab})
  =
  T(h_b\circ h_a)
  =
  T(h_b)\circ T(h_a)
  \qquad\text{et}\qquad
  T(h_{a+b})
  =
  T(h_a+h_b)
  =
  T(h_a)+T(h_b);
\]
cela prouve que l'application $(a,y)\to T(h_a)(y)$
est une loi externe de $A$-module à droite sur $T(M)$. En particulier,
$T(A_s)$ est un $A$-\emph{module à droite}.
\end{env}

\begin{env}[7.1.2]
Lorsque $A$ est un anneau \emph{commutatif}, il résulte de (7.1.1) que pour tout
$A$-module $M$, $T(M)$ est naturellement muni d'une structure de $A$-module;
en outre, si $u:M\to N$ est un homomorphisme de $A$-modules, on a, pour tout
$a\in A$, $u\circ h_{a,M}=h_{a,N}\circ u$,
d'où $T(u)\circ T(h_{a,M})=T(h_{a,N})\circ T(u)$,
ce qui prouve que $T(u):T(M)\to T(N)$ est un homomorphisme de $A$-\emph{modules}; on
voit donc que $T$ peut être considéré comme un foncteur covariant additif de
la catégorie $\boldsymbol{Ab}_A$ dans elle-même. De façon précise, on a ainsi
défini une \emph{équivalence} canonique entre la catégorie des foncteurs covariants
additifs $\boldsymbol{Ab}_A\to\boldsymbol{Ab}$ et la catégorie des foncteurs
covariants $A$-\emph{linéaires}
$T:\boldsymbol{Ab}_A\longrightarrow\boldsymbol{Ab}_A$,
c'est-à-dire tels que
$T(h_{a,M})=h_{a,T(M)}$ pour tout $a\in A$.
Comme le foncteur d'inclusion $I:\boldsymbol{Ab}_A\longrightarrow\boldsymbol{Ab}$,
faisant correspondre à tout $A$-module le $\mathbf Z$-module sous-jacent, est
exact et fidèle, les propriétés d'exactitude des deux foncteurs associés par
l'équivalence précédente sont les mêmes.
\end{env}

\begin{env}[7.1.3]
L'anneau $A$ étant toujours supposé commutatif, soit $B$ une $A$-\emph{algèbre} (non
nécessairement commutative), et soit $\rho:A\to B$ l'homomorphisme d'anneaux
correspondant à cette structure d'algèbre; cet homomorphisme définit un
foncteur covariant
